\section{Calculi from Certain States}
\label{sec:calculi}

We describe how to construct, for a given set of
events together with a set of assignments that settle every event, 
the set of mixtures of these assignments. The construction and its
completeness theorem are de Finetti's~\cite{definetti1974}, in the form
that Paris and Williams use for non-classical
logics~\cite{paris2001,williams2012}. All functions take values in the
nonnegative reals $\mathbb{R}_{\ge 0}$.\footnote{The Lean development
works in the extended nonnegative reals $[0,\infty]$ of \textsf{mathlib},
the type in which measures take values. Every function in this paper takes
finite values, so the statements are the same, except that some Lean statements
carry a finiteness hypothesis that is automatic over $\mathbb{R}_{\ge 0}$,
as in Theorem~\ref{thm:tensor-collapse}.}

\subsection{Presentations and Mixtures}
\label{sec:calculi-mix}

\begin{definition}[presentation, calculus\leanref{mix}]
\label{def:mix}
A \emph{presentation} is a set $E$, whose elements are the \emph{events},
together with a set $S$ of functions $E \to \mathbb{R}_{\ge 0}$ with values in
$\{0,1\}$, the \emph{certain states}. A \emph{mixture} of functions
$s_1, \dots, s_n$ is a function $\sum_{i} w_i\, s_i$ with weights
$w_i \in \mathbb{R}_{\ge 0}$ and $\sum_i w_i = 1$. The \emph{calculus} $\Mix(S)$ is
the set of mixtures of certain states. A set of functions is
\emph{closed under mixtures} if it contains every mixture of its
elements.
\end{definition}

\noindent
Mixtures are finite. When $E$ is finite, so is $S$, and a countable
mixture of elements of $\Mix(S)$ regroups into a finite mixture of
certain states, so countable mixtures would change nothing. When $E$ is
infinite they differ. If $S$ contains the indicators of the singletons
$\{e_1\}, \{e_2\}, \dots$, then $\sum_n 2^{-n}\, 1_{\{e_n\}}$ takes
infinitely many values, while a finite mixture of $\{0,1\}$-valued
functions takes finitely many, so it is not in $\Mix(S)$.

\noindent
A certain state settles every event, with value $1$ or $0$. A mixture
records uncertainty about which certain state obtains, with $w_i$ the
weight of $s_i$. In Section~\ref{sec:linear} the events are the points of
the web of a formula of linear logic and the certain states are the
deterministic behaviours of that type, so an element of the calculus is a
probability distribution over deterministic behaviours.

\begin{theorem}[certain elements and mixing\leanref{sharp-mem-mix-iff}\leanref{mix-mix}]
\label{thm:mix-basic}
For every presentation $(E, S)$:
(i) a function with values in $\{0,1\}$ lies in $\Mix(S)$ iff it lies in
$S$, and
(ii) $\Mix(\Mix(S)) = \Mix(S)$, so $\Mix(S)$ is the least set that
contains $S$ and is closed under mixtures.
\end{theorem}

\noindent
For (i), let $x = \sum_i w_i s_i$ have values in $\{0,1\}$, and let
$w_j \ne 0$. At an event $e$ with $x_e = 0$, every $s_i$ with
$w_i \ne 0$ has $s_i(e) = 0$. At an event with $x_e = 1$, every such
$s_i$ has $s_i(e) = 1$, since $\sum_i w_i = 1$ and $s_i(e) \le 1$. So
$s_j = x$. Part (ii) expands a mixture of mixtures into one mixture with
product weights. Mixing therefore creates no new certain states, and the
certain states of the calculus are those of the presentation.

\subsection{Affine Laws and Completeness}
\label{sec:calculi-laws}

The calculus has an axiomatic description by linear inequalities.

\begin{definition}[affine law\leanref{affinelaw}]
\label{def:affine-law}
Let $E$ be finite. An \emph{affine law} on $E$ is given by functions
$a, b : E \to \mathbb{R}_{\ge 0}$ and constants $k, k' \in \mathbb{R}_{\ge 0}$. A function
$x : E \to \mathbb{R}_{\ge 0}$ \emph{satisfies} it if
\[
  \sum_{e \in E} a_e x_e + k \;\le\; \sum_{e \in E} b_e x_e + k'.
\]
An affine law is \emph{valid on} a set of functions if every function in
the set satisfies it.
\end{definition}

\noindent
Every linear inequality with real coefficients between the values of
$x$ can be written in this form, by moving negative coefficients to the
other side. For example, $x_e \le 1$, $x_e + x_{e'} \le 1$, and the
equation $x_e + x_{e'} = x_{e''}$, as two inequalities, are affine laws.

\begin{theorem}[completeness~\cite{definetti1974,paris2001}\leanref{mix-eq-laws}\leanref{affinelaw-holds-of-mix}]
\label{thm:mix-eq-laws}
Let $E$ be finite and $S$ a set of functions $E \to \mathbb{R}_{\ge 0}$ with values
in $\{0,1\}$. Then $\Mix(S)$ is the set of functions $x : E \to \mathbb{R}_{\ge 0}$
that satisfy every affine law valid on $S$.
\end{theorem}

\noindent
Affine laws are preserved by mixtures, which gives the inclusion from
left to right. For the converse, a point outside the convex hull of the
finite set $S$ is separated from it by a hyperplane, and the hyperplane
is an affine law valid on $S$ that the point violates. The Lean proof uses
the geometric Hahn--Banach theorem of \textsf{mathlib}. For rational
points and certain states, the development also contains a proof without
the axiom of choice by Fourier--Motzkin elimination, the elimination of
variables from a system of linear inequalities\leanref{constr-farkas}.

Read as bets, an affine law valid on $S$ and violated at $x$ is a finite
system of bets, priced by $x$, that loses money in every certain state.
Paris and Williams use Theorem~\ref{thm:mix-eq-laws} in this form for
non-classical logics, and Williams derives from it that the mixtures are
the assignments that no other assignment dominates in
accuracy~\cite{paris2001,williams2012}.

\subsection{Laws of Deterministic Programs}
\label{sec:calculi-unique}

The calculus is determined by its certain states and their laws.

\begin{theorem}[uniqueness\leanref{calculus-unique}\leanref{calculus-laws-iff}]
\label{thm:unique}
Let $E$ be finite, let $S$ be a set of functions $E \to \mathbb{R}_{\ge 0}$ with
values in $\{0,1\}$, and let $C$ be a set of functions
$E \to \mathbb{R}_{\ge 0}$ that contains $S$.
(i) $C$ satisfies no affine law that some element of $S$ violates. If
$C$ satisfies every affine law valid on $S$, then the affine laws valid
on $C$ are exactly those valid on $S$.
(ii) If $C$ is closed under mixtures and satisfies every affine law valid
on $S$, then $C = \Mix(S)$.
\end{theorem}

\noindent
Part (i) holds because $S \subseteq C$. For (ii), closure under mixtures
gives $\Mix(S) \subseteq C$, and Theorem~\ref{thm:mix-eq-laws} gives
$C \subseteq \Mix(S)$. So $\Mix(S)$ is at once the least set that
contains the certain states and is closed under mixtures and the largest
set that obeys their laws. The mathematics is de Finetti's and
Paris's~\cite{definetti1974,paris2001}.

In the rest of the paper $C$ is the probabilistic coherence space of a
formula, which contains its certain states and is closed under mixtures,
and $\Mix(S)$ is the calculus of the same formula. The question is which
of the two cases of the following definition holds.

\begin{definition}[gap]
\label{def:gap}
Let $(E, S)$ be a presentation and $C \supseteq S$ a set of functions
$E \to \mathbb{R}_{\ge 0}$ that is closed under mixtures. Then $\Mix(S) \subseteq C$,
and $C$ has a \emph{gap} over $S$ if the inclusion is strict.
\end{definition}

\begin{corollary}[a gap is a violated law\leanref{mix-eq-laws}]
\label{cor:gap-law}
For finite $E$, the set $C$ of Definition~\ref{def:gap} has a gap over
$S$ iff some element of $C$ violates an affine law valid on $S$.
\end{corollary}

\noindent
This is Theorem~\ref{thm:mix-eq-laws} applied to the elements of $C$
outside $\Mix(S)$. When the certain states are the deterministic
programs of a type, an affine law valid on $S$ is an inequality that
every deterministic program obeys, and every probabilistic program that
flips coins and then runs a deterministic program obeys it too. A gap is
an element of the probabilistic coherence space that violates such an
inequality. The witnesses found below are of this kind: a five-term
inequality at the type of a program with two functional arguments, with
the bounds of the KCBS inequality (Section~\ref{sec:second}), and the
CHSH inequality at the dual type (Section~\ref{sec:duality}).

\subsection{Why Not Lattice Valuations}
\label{sec:calculi-lattice}

For intuitionistic and classical logic, the calculus of a finite event
algebra is also given by the familiar axioms of a valuation: a grading
$v$ of the events with $v(\bot) = 0$ and $v(\top) = 1$ that is monotone
and \emph{modular}, $v(a) + v(b) = v(a \sqcup b) + v(a \sqcap b)$, where
$\sqcup$ and $\sqcap$ are the join and meet of the algebra of
events~\cite{paris2001}\leanref{frame-recipe}. The event algebras of linear logic are
\emph{quantales}, complete lattices with an associative multiplication
$\otimes$ that distributes over arbitrary joins and need not be
idempotent~\cite{girard1987}. The tensor $\otimes$ is not the meet, so
modularity does not apply to it as it stands. Replacing the meet by the
tensor in the modularity axiom does not work either.

\begin{definition}[tensor valuation\leanref{tensor-valuation}]
\label{def:tensor-valuation}
Let $\Omega$ be a \emph{join-semilattice}, a partial order in which every
two elements have a least upper bound $a \sqcup b$, and let $\otimes$ be
a binary operation on $\Omega$. A \emph{tensor valuation} for $\otimes$
is a function $v : \Omega \to \mathbb{R}_{\ge 0}$ such that
$v(a) + v(b) = v(a \sqcup b) + v(a \otimes b)$ for all $a, b \in \Omega$.
\end{definition}

\begin{theorem}[tensor collapse\leanref{tensor-collapse}\leanref{tensor-dichotomy}]
\label{thm:tensor-collapse}
Let $\Omega$ be a join-semilattice, $\otimes$ a binary operation on
$\Omega$, and $v$ a tensor valuation for $\otimes$. Then
$v(a \otimes a) = v(a)$ for every $a \in \Omega$.
\end{theorem}

\noindent
Setting $b = a$ in the defining equation and using $a \sqcup a = a$ gives
$v(a) + v(a) = v(a) + v(a \otimes a)$, and $v(a)$ cancels. No property of $\otimes$ is used.

\begin{example}[a constant valuation\leanref{tensor-witness}]
\label{ex:tensor-witness}
On $\mathbb{N}$ with $\sqcup = \max$ and $\otimes = +$, an associative
operation that distributes over $\max$ and has $1 + 1 \ne 1$, setting $a = 0$ in the defining equation gives
$v(n) = v(0)$ for every $n$, so every tensor valuation is constant.
\end{example}

\noindent
So a grading of the events of a quantale by this law cannot separate one
copy of a resource, $a$, from two copies, $a \otimes a$, which is the
distinction that the tensor of linear logic exists to make.
Probabilistic coherence spaces make the distinction by assigning weights
to the points of a web, with $a \otimes a$ interpreted on a product of
webs, and not by grading formulas~\cite{danos-ehrhard2011}. The
construction of this section applies to them, because it needs a set of
events and a set of certain states and no operations on events.
